\documentclass[a4paper]{article} %

\date{} %

% Please, do not change any of the above lines

\usepackage{amssymb} % add other LaTeX packages you need here.

\begin{document}
\setcounter{page}{1} % Please, do not delete this line

% Your title
\title{Stability of Mapping Spaces in $\varepsilon$-Simplicial Sets}

% Authors
\author{Dominik Lachman\\ %
       Palacký University Olomouc\\ \\ % Affiliation 1
       % New author \\
       % New affiliation \\
       % Add authors and affiliation as needed
       \texttt{dominik.lachman@upol.cz} % Only one corresponding e-mail, please
       }%

\maketitle

% The abstract

A familiar stability phenomenon in simplicial homotopy theory is that horn-filling conditions characterizing categories and groupoids are preserved under mapping-space constructions. For instance, simplicial sets satisfying the inner horn-filling conditions retain this property after passing to suitable mapping spaces, and an analogous statement holds for Kan complexes.


In this talk, I will discuss a related phenomenon for $\varepsilon$-simplicial sets. These were introduced in the simplicial study of Frobenius monoids in the category $\mathbf{Rel}$ of sets and relations \cite{Lachman2026}, \cite{Lachman2025}. Within this framework, we consider a distinguished class of $\varepsilon$-fibrant objects, characterized by appropriate filling conditions.

A key ingredient is the close relationship between $\varepsilon$-simplicial sets and paracyclic sets. This connection provides a useful way to reinterpret certain simplicial filling problems and plays an essential role in the analysis of $\varepsilon$-fibrancy.

The main result is a stability theorem showing that, under suitable hypotheses, the relevant $\varepsilon$-fibrancy conditions are preserved under mapping-space constructions. I will explain the analogy with the classical cases of categories and groupoids, the role of paracyclic structures, and the main ideas of the proof.

\begin{thebibliography}{99}

\bibitem{Lachman2026}
D.~Lachman.
\newblock \emph{Simplicial approach to Frobenius algebras in the category of relations,}
\newblock Appl. Categ. Structures \textbf{34}, 25 (2026).

\bibitem{Lachman2025}
D.~Lachman.
\newblock \emph{Simplicial perspectives on (pseudo) effect algebras,}
\newblock arXiv:2510.02748 (2025).

\end{thebibliography}

\end{document}